\documentclass[12pt]{article}
\usepackage{graphicx}
\usepackage{geometry}
\usepackage{listings}
\usepackage[T1]{fontenc}
\usepackage{url}
\usepackage[hidelinks]{hyperref}
\usepackage{float}
\usepackage{caption}
\usepackage{bookmark}
\geometry{margin=1in}

\lstset{
  basicstyle=\ttfamily\small,
  breaklines=true,
  breakatwhitespace=true,
  columns=fullflexible
}

% ---- Macro for bottom-only cropping ----
% Usage: \cropbottom[<bottom>]{<file>}
% Example: \cropbottom[150pt]{step0.png}
\newcommand{\cropbottom}[2][150pt]{%
  \includegraphics[width=0.9\textwidth, trim={0 #1 0 0}, clip]{#2}%
}

\begin{document}

\begin{center}
    {\Large \textbf{COMP 324 Network Security}} \\[4pt]
    {\large Module 3 Lab 1: Linux User \& Group Management} \\[10pt]
    \textbf{Lab Report} \\[6pt]
    \textbf{Student: Erik Gonzalez} \\[4pt]
    \textbf{Professor: Samuel Addington} \\[4pt]
    \textbf{Date: \today}
\end{center}

\section*{Exact Commands Run}
\begin{lstlisting}
# Step 0: Verify sudo
sudo -v

# Step 1: Add new user student1
sudo useradd -m -s /bin/bash student1
sudo passwd student1
id student1
getent passwd student1

# Step 2: Create group developers
sudo groupadd developers
getent group developers

# Step 3: Add student1 to developers
sudo usermod -a -G developers student1
id student1
su - student1
groups
exit
\end{lstlisting}


\section*{Step 0: Verify Sudo Access}
\begin{figure}[H]
    \centering
    \cropbottom[1200pt]{step0.png}
    \end{figure}



\section*{Step 1: Add New User \texttt{student1}}
\begin{figure}[H]
    \centering
    \cropbottom[800pt]{step1.png}
    
\end{figure}



\section*{Step 2: Create Group \texttt{developers}}
\begin{figure}[H]
    \centering
    \cropbottom[1100pt]{step2.png}
    
\end{figure}



\section*{Step 3: Add \texttt{student1} to \texttt{developers}}
\begin{figure}[H]
    \centering
    \cropbottom[900pt]{step3.png}
    
\end{figure}



\section*{Reflection}

I found it surprising how Linux separates users into different groups, each with its own permissions, and how this affects what commands and actions a user can perform.  




\end{document}

